\subsection{Límites}

\begin{itemize}
  \item El videojuego no será un laboratorio técnico de ciberseguridad ni un entorno de virtualización real; por tanto, no incluirá máquinas virtuales, simulación de redes reales, explotación de vulnerabilidades reales ni ejecución de ataques sobre sistemas verdaderos.
  \item El contenido del proyecto se limitará a un conjunto seleccionado de amenazas digitales fundamentales, sin abarcar todos los dominios avanzados de la ciberseguridad, como criptografía profunda, análisis forense digital, pentesting avanzado o respuesta profesional a incidentes.
  \item El prototipo será monojugador y de ejecución local, por lo que no contemplará arquitectura multijugador, sincronización en red, servicios en línea, tablas de clasificación globales ni interacción simultánea entre usuarios.
  \item La herramienta será un apoyo pedagógico complementario y no reemplazará la labor del docente, la currícula académica, los laboratorios especializados ni la asesoría profesional necesaria para incidentes reales de seguridad.
  \item La validación del proyecto se enfocará en la funcionalidad del prototipo y en la usabilidad de la experiencia de aprendizaje, sin incluir un estudio estadístico de largo plazo sobre retención de conocimiento o impacto pedagógico posterior al uso.
  \item El módulo de seguimiento registrará información interna de desempeño dentro del videojuego, pero no se integrará con sistemas institucionales de calificación, plataformas académicas, bases de datos externas ni servicios de gestión educativa.
\end{itemize}
